% Generated draft: requires mathematical and source review.
\documentclass[11pt]{article}
\usepackage[margin=25mm]{geometry}
\usepackage{fontspec}
\setmainfont{Latin Modern Roman}
\usepackage{amsmath,amssymb,mathtools,unicode-math}
\setmathfont{Latin Modern Math}
\usepackage{enumitem,needspace,xurl,hyperref}
\setlength{\parindent}{0pt}
\setlength{\parskip}{5pt}
\setlength{\emergencystretch}{3em}
\newcommand{\N}{\mathbb{N}}
\newcommand{\Z}{\mathbb{Z}}
\newcommand{\Q}{\mathbb{Q}}
\newcommand{\R}{\mathbb{R}}
\newcommand{\C}{\mathbb{C}}
\newcommand{\F}{\mathbb{F}}
\newcommand{\A}{\mathbb{A}}
\newcommand{\PP}{\mathbb{P}}
\newcommand{\E}{\mathbb{E}}
\newcommand{\eps}{\varepsilon}
\newcommand{\dd}{\,\mathrm{d}}
\begin{document}
\section*{UnsolvedMath ID: 30002101}
\textbf{Source problem code:} catalogue entry
\subsection{Definitions and mathematical statement}
The symbol \(\overline{\mathcal M}_g\) conventionally denotes the moduli stack of stable curves of arithmetic genus \(g\). In the usual unmarked setting, a stable curve is a proper, geometrically connected nodal curve with finite automorphism group. The source uses \(\overline{\mathcal M}_g(\alpha)\) for a log canonical model associated with this moduli problem, with parameter \(\alpha\). It does not specify the ground field, the permitted genera or parameter values, the defining log canonical divisor, or whether this notation denotes a stack or its associated coarse moduli space. These choices remain unresolved here.

\textbf{Question.} Does the log canonical model
\[
\overline{\mathcal M}_g(\alpha)
\]
inherit a moduli interpretation as a moduli space of singular curves?

Such an interpretation would require specifying the curves allowed, their stability conditions, their families and isomorphisms, and the relationship between their moduli object and the indicated log canonical model. The supplied question does not specify these data or the required sense of moduli space, so it does not determine a precise representability or coarse-moduli assertion. The accompanying assessment describes a uniform question in genus and \(\alpha\), but supplies no exact quantifier ranges.

\textbf{Supplied qualification.} The record labels the problem partially solved. Its dated assessment states that several early log canonical models have modular descriptions through pseudostable and related singular curves, that the genus-four program has recently been completed, and that a uniform modular description for all parameters and genera is not known. These are supplied curation claims, not independently verified mathematical conclusions of this draft.
\subsection{Short English statement}
The question asks whether log canonical models of the moduli of stable curves can themselves classify suitable singular curves. The allowed genera, parameters, ground field, singularities, stability conditions, and meaning of moduli space remain unspecified. The supplied partially solved status and genus-four completion claim require independent review.
\subsection{Sources}
\begin{itemize}
\item Local UnsolvedMath registry, supplied source record with permanent ID \(30002101\), source problem code \texttt{catalogue entry}, in the collection \emph{Oberwolfach Reports Open Problems}. The project identifies \texttt{data/unsolvedmath/problems.json} and the supplementary \texttt{catalog/problems.json} as registry files; the supplied record does not identify which file contains this entry. No individual proposer is supplied. The proposed year is recorded as 2012.
\item Supplied report citation: \emph{Classical Algebraic Geometry (2013)}, Oberwolfach Reports, pp.\ 1845--1893, DOI \texttt{10.4171/owr/2012/30}. \url{https://doi.org/10.4171/owr/2012/30}. The background describes Oberwolfach Reports as documenting mathematical workshops and their open-problem discussions. The citation's 2013 designation and the record's 2012 dating are retained without reconciliation.
\item Supplied literature reference: Hassett and Hyeon, first flip. \url{https://arxiv.org/abs/0806.3444}.
\item Supplied literature reference: Hassett--Keel program in genus four. \url{https://arxiv.org/abs/2504.20312}.
\item Supplied literature reference: original OWR report. \url{https://doi.org/10.4171/owr/2012/30}.
\item Literature assessment dated exactly \texttt{2026-08-21}: supplied dated curation, not independently verified. It labels the problem partially solved and makes the early-model, genus-four completion, and remaining uniform-description claims reproduced above. The supplied background explicitly requests independent verification of time-sensitive status and cited claims.
\item The registry labels the statement \texttt{corrected\_verified} and reports that its mathematical content was checked against the official Oberwolfach report, with notation, context, and question form normalized without changing the source-level claim. This is a supplied verification statement; this draft does not independently verify it or claim to have consulted the report.
\end{itemize}
\end{document}
